\section{Prueba de entrada}

\begin{frame}
	\frametitle{\insertsection}

	\begin{enumerate}
		\item

		      Una barra metálica de \qty{100}{\centi\metre} de longitud tiene sus
		      extremos $x=0$, $x=100$ mantenidos a \qty{0}{\degreeCelsius}.
		      Inicialmente, la mitad de la barra está a \qty{60}{\degreeCelsius};
		      mientras que la otra mitad a \qty{40}{\degreeCelsius}.
		      Asumiendo un coeficiente de difusividad de \num{0.16} unidades c.g.s
		      y un entorno aislado.

		      \begin{enumerate}
			      \item

			            Describa la ecuación diferencial correspondiente, así como las
			            condiciones iniciales y de frontera.

			      \item

			            No resuelva, solo comente los pasos a tener en cuenta para resolver
			            el problema en el ítem~(a).
		      \end{enumerate}

		\item

		      Dado el problema
		      \begin{equation*}
			      \begin{cases}\displaystyle
				      \diffp{u}{t}=
				      \diffp[2]{u}{x}+
				      a\diffp{u}{x} &
				      \text{en }\mathbb{R}\times\left(0,\infty\right). \\
				      u=f           &
				      \text{en }\mathbb{R}\times\left\{t=0\right\}.
			      \end{cases}
		      \end{equation*}
		      Utilizando la transformada de Fourier, determine la solución en el contexto de su transformada.

		\item

		      \begin{enumerate}
			      \item

			            Describa cuatro términos de la serie de Taylor de una función real
			            de variable real infinitamente diferenciable en un dominio
			            adecuado.

			      \item

			            Describa cuatro términos de la serie de una función real de dos
			            variables reales infinitamente diferenciable en un dominio
			            adecuado.

			      \item

			            Describa el gradiente y el hessiano de una función real de tres
			            variables reales en un dominio adecuado.
		      \end{enumerate}
	\end{enumerate}
\end{frame}
